\section{Discussion}\label{sec:discussion}
Material: per-frame scale analysis on 6 scenes (\code{analyze\_scale\_drift.py}, stride 5, 471 frames), the failure-frame
study of scene 42444474, and Tables~\ref{tab:main}--\ref{tab:exp1scene}. Per-frame AbsRel uses protocol~2.

\subsection{Why the LiDAR-teacher fine-tune works---and why it is not enough}\label{sec:whyft}
We ran the per-frame analysis of \S\ref{sec:metric} on R3 as well (every fifth frame of the six test rooms, 471
frames): per frame, the GT/pred ratio (median over the frame) and AbsRel under different amounts of scale help
(Table~\ref{tab:whyft}).
\begin{table*}[t]
\centering\small
\caption{Per-frame scale of the pretrained metric model vs.\ the LiDAR-teacher fine-tune R3, six test rooms, every fifth
frame. AbsRel uses protocol~2.}
\label{tab:whyft}
\begin{tabularx}{\textwidth}{@{}lllX@{}}
\toprule
 & Pretrained & \textbf{R3} & Meaning \\
\midrule
median(GT / pred) per room & 0.69--0.85 & \textbf{0.93--1.21} & \textbf{the domain bias is gone}: pretrained is 15--31\,\% too far in every room; R3 is within $\sim$20\,\% and not always in the same direction \\
per-frame swing within a room (p95/p5 of GT/pred) & 1.33--1.79$\times$ (mean 1.55) & 1.26--1.92$\times$ (mean 1.52) & \textbf{fine-tuning does not reduce the per-frame swing} \\
AbsRel, raw & 0.381 & \textbf{0.134} & \\
AbsRel with one $(S,T)$ per room fitted on 10 GT frames & 0.127 & 0.117 & pretrained $+$ per-room GT calibration $\approx$ R3 with no calibration at all \\
AbsRel with the right scale on every frame (one factor) & 0.045 & 0.051 & the per-frame scale is almost all of R3's remaining error \\
AbsRel, oracle $(s,t)$ per frame & 0.035 & 0.042 & shape hardly changes (frozen encoder) \\
\bottomrule
\end{tabularx}
\end{table*}
\begin{enumerate}
\item \textbf{Fine-tuning works because it learns the domain's multiplier.} R3 with no calibration (AbsRel 0.134) is on
par with the pretrained model given a per-room scale fitted on ground truth (0.127): what a one-off calibration per room
provides has been moved into the head's weights, using LiDAR labels from other rooms. With the encoder frozen, shape
barely changes (oracle 0.035 $\rightarrow$ 0.042); the whole gain is scale.
\item \textbf{It is not enough because the per-frame swing remains.} Within one room, frames at p5 and p95 still need
factors $\sim$1.5$\times$ apart, as before fine-tuning, and that swing is nearly all of what is left: with the right
scale on every frame R3's AbsRel drops from 0.134 to 0.051 (90\,\% of the way to the oracle), whereas one $(S,T)$ per
room only reaches 0.117. The failing frames are those of \S\ref{sec:drift}--\ref{sec:metric} (blank walls, white
cabinets, close downward views with no size cue). The missing information is not in a single image and does not come
from a better teacher (the laser teacher does no better, \S\ref{sec:finetune}); it has to come from metric data at run
time.
\item \textbf{Fine-tuning and per-frame metric points are complementary, not alternatives.} Fine-tuning removes the
domain bias with nothing at run time (\SI{\rsMETRICCM}{cm} $\rightarrow$ \SI{\rsFTLIDARALLCM}{cm}); per-frame sparse
points remove the per-frame swing (relative model: \SI{\rsSCENECM}{cm} $\rightarrow$ \SI{\rsSPARSECM}{cm},
\S\ref{sec:exp1}). The 2D bound for combining them is the 0.051 ``right scale on every frame'' row above; measuring
R3 $+$ \code{sparse\_points} in 3D is the most direct next step (\S\ref{sec:limitations}).
\end{enumerate}

\subsection{Why per\_scene loses to the oracle by $4.5\times$: the model's scale follows what is in the frame, not time}\label{sec:drift}
An affine-invariant model promises $\text{disparity} \approx a\cdot(1/\text{depth}) + b$ with $(a,b)$ depending on
the image. If $(a,b)$ were constant over a scan, one calibration (\code{per\_scene}) would equal the oracle. We measured
how constant it is by refitting $(s,t)$ on every fifth frame (Fig.~\ref{fig:drift}, Table~\ref{tab:drift}).
\begin{figure*}[t]
\centering
\includegraphics[width=\textwidth]{scale_drift_timeline}
\caption{Per-frame oracle scale divided by the per-scene scale, over time, for DA-v2 Large on all six scenes (log axis;
1.0 $=$ the one-off calibration was right for that frame). Orange dots mark close-up frames (median GT depth
$<$\SI{1}{m}). The ratio swings by $2$--$4\times$ within every scan with no monotone trend.}
\label{fig:drift}
\end{figure*}
\begin{table*}[t]
\centering\small
\caption{DA-v2 Large, per-scene values across the six scenes.}
\label{tab:drift}
\begin{tabularx}{\textwidth}{@{}lXX@{}}
\toprule
 & Range over scenes & Meaning \\
\midrule
$s$ p95 / p5 within one scene & \textbf{2.6--4.9$\times$} & the central 90\% of frames need scales 3--5$\times$ apart \\
$s$ p75 / p25 & 1.47--1.56 & even the middle half swings $\pm25\%$ \\
corr($s$, time) & $-0.33 \ldots +0.30$ (sign flips) & \textbf{not cumulative drift}: no direction \\
corr($s$, frame's median depth) & \textbf{$-0.27 \ldots -0.75$ (negative in every scene)} & close-up frames need a different $s$: view-dependent \\
AbsRel oracle $\rightarrow$ per\_scene & 0.032 $\rightarrow$ 0.239 (\textbf{7.5$\times$}) & in 2D, before fusion, one scale costs $7\times$ \\
per\_scene without close-ups ($<$\SI{1}{m}) & 0.152 (still $5\times$ oracle) & close-ups hurt but are not the only cause \\
scale-only per frame, depth space & 0.331 (worse than per\_scene) & a disparity shift is needed in every scene \\
\bottomrule
\end{tabularx}
\end{table*}
\begin{figure}[t]
\centering
\includegraphics[width=\columnwidth]{scale_vs_depth}
\caption{The same ratio against the frame's median GT depth, all scenes pooled: the scale a frame needs correlates with
how close the camera is to what it sees.}
\label{fig:scalevsdepth}
\end{figure}
Three consequences:
\begin{enumerate}
\item \textbf{No number of fit frames helps \code{per\_scene}}: the $(S,T)$ fitted on 10 spread frames already lands
within $\pm13\%$ of the median oracle scale in five of the six scenes (0.87--1.17$\times$; the exception, 47333774,
is 0.67$\times$). The problem is that the median does not apply to most frames, not that it is estimated badly.
\item \textbf{A $7.5\times$ per-frame error becomes a $4.5\times$ 3D error} not because TSDF averages it out but because
the error concentrates in a minority of frames: frames with per\_scene AbsRel $>1$ (placed $\geq2\times$ wrong) are
downward shots of table-tops or floor at 0.5--\SI{0.8}{m} and frame-filling blank walls or cabinets. These frames carry
no size cue, so the model guesses disparity at ``room scale''; multiplied by the room's $(S,T)$, a table at \SI{0.5}{m}
lands at 1.5--\SI{2}{m} as a phantom surface in the middle of the room---the red blob in every per\_scene panel of
Fig.~\ref{fig:exp1}.
\item \textbf{The effect is general, not a property of the development scene}: 47429736 (bedroom$+$closets) has
close-ups in 53/101 frames yet per\_scene AbsRel 0.197, the same as a scene with 3/93 (45261556: 0.172); and the five
unseen scenes give per\_scene Chamfer 19.9--\SI{26.0}{cm}, bracketing 42444474 (24.1).
\end{enumerate}
\textbf{System implication.} The assumption ``one scale per scene'' is false for this model family on real room video.
Two ways remain: (a) \emph{per-frame} metric data---3D feature points from VIO, which ARKit/ARCore deliver with the pose
anyway (\code{sparse\_points}); or (b) drop or down-weight cue-less frames before fusion---cheaper, but removing
close-ups still leaves $4\times$ the oracle, so it can only supplement (a). Exp.~1 confirms (a): 200 points per frame
bring the row to \SI{\rsSPARSECM}{cm}, \SI{\rsSPARSEGAP}{cm} from the oracle, with the same per-scene ordering. The per-frame
fit needs no history and tolerates frames without points by carrying the previous $(s,t)$, exactly what a tracker
outage looks like.

\subsection{Why the ``metric'' model gives \SI{\rsMETRICCM}{cm}: right on average, wrong per frame}\label{sec:metric}
DA-v2 Metric-Indoor should output metres directly. Per-frame analysis over the six scenes:
\begin{table*}[t]
\centering\small
\begin{tabularx}{\textwidth}{@{}lXX@{}}
\toprule
 & Range over scenes & Meaning \\
\midrule
median(GT / pred) per scene & 0.69--0.85 & the model says 15--31\% too far on average (per-scene bias) \\
GT / pred within a scene & e.g.\ 0.41--0.86, 0.43--1.35 & frames in one room are wrong from $-59\%$ to $+35\%$ \\
oracle $t$ in inverse space & $-0.11 \ldots -0.02$ ($\approx 0$) & \textbf{no shift}: the shape is right in the affine sense \\
AbsRel raw $\rightarrow$ scale-only/frame $\rightarrow$ oracle $(s,t)$ & 0.381 $\rightarrow$ \textbf{0.045} $\rightarrow$ 0.035 & one scalar per frame recovers 97\% of the gap \\
AbsRel with one $(S,T)$ per scene & 0.127 & like per\_scene for the relative model: the median is not enough---and what fine-tuning moves into the weights (\S\ref{sec:whyft}) \\
\bottomrule
\end{tabularx}
\end{table*}
In 3D the same surface from several frames is placed at several distances and TSDF averages them into a thick slab:
accuracy \SI{\rsMETRICACC}{cm} but completeness \SI{\rsMETRICCOMP}{cm} (Table~\ref{tab:exp1})---everything is there,
in the wrong place. The worst scene (kitchen 45261556, \SI{\rsMETRICf}{cm}) is among the model's best shapes under the
oracle (\SI{\rsORACLEf}{cm}), confirming a pure scale effect. The worst frames in every scene look alike: a single flat, uniformly coloured surface
filling the frame (white cabinets, tiles) with nothing distant to compare against---the model's prior pushes flat
surfaces away. \textbf{Implication:} a metric model does not remove calibration; it reduces it from two parameters per
frame $(s,t)$ to \textbf{one} $(s)$, which is a real advantage: a handful of VIO points suffice to fit one scalar
stably, whereas $(s,t)$ needs points spread over the depth range. The other way to get that scalar is to put it in
the weights: fine-tuning the same checkpoint with LiDAR labels from real iPad frames (R3, \num{19896} frames) moves the
row to \SI{\rsFTLIDARALLCM}{cm} with no run-time fit at all (\S\ref{sec:main}), the deployable answer for a device that
has neither LiDAR nor reliable VIO points; \S\ref{sec:whyft} shows that what it fixes is the domain bias, while the
per-frame swing remains.

\subsection{Which rooms are hard, and for whom}\label{sec:rooms}
Ranking scenes by Chamfer per row (Table~\ref{tab:exp1scene}): Faro/LiDAR go from bedroom 47333774 (easiest) to the
bathroom (\rsGTb--\rsGTd{} / \rsLIDARb--\rsLIDARd); mono$+$oracle from the development kitchen 42444474 and kitchen
45261556 to the bathroom (\rsORACLEa--\rsORACLEd); mono$+$per\_scene from living room 47115299 to bedroom 47333774
(\rsSCENEc--\rsSCENEb).
\begin{itemize}
\item \textbf{The bathroom is hard for every source} (LiDAR $+1.3$ from Faro, oracle $+3.4$ from LiDAR): the mirror yields
phantom depth for LiDAR and model alike, glossy tiles reflect; it is the only scene where the oracle fit produced
$s \leq 0$ on some frames (a frame-filling mirror). If the target market is bathroom renovation, this paper's averages
are optimistic.
\item \textbf{Standard bedroom 47333774 is easiest for LiDAR but second-hardest for mono} (oracle \rsORACLEb): a small room
scanned close to walls and bed leaves most frames at 1--\SI{2}{m} of blank wall with no cue. Rooms with ``lots of
stuff'' (kitchen 45261556: 10 cabinets; living room: 4 tables, 2 sofas) are easy for mono---the model likes objects of
known size.
\item The mono difficulty order does not match LiDAR's $\Rightarrow$ \textbf{``hard scenes'' must be assessed per
source}; LiDAR experience does not predict where mono fails.
\end{itemize}

\subsection{Visible only in 3D: mono disagrees with itself across views}
Exp.~3 shows what 2D metrics cannot: reducing frames from 391 to 20 (stride 20) \textbf{improves mono$+$oracle accuracy
from \rsSTRORACLEONEACC{} to \SI{\rsSTRORACLETWENTYACC}{cm}} while 2D AbsRel stays at
\rsSTRORACLEONEABSREL--\rsSTRORACLETWENTYABSREL{} and GT accuracy at \rsSTRGTONEACC--\SI{\rsSTRGTTWENTYACC}{cm}. Even with every frame perfectly scaled,
the shape of the same surface from different views \textbf{disagrees} at the $\sim$\SI{2}{cm} level, and fusing more
frames thickens the surface rather than sharpening it---the multi-view consistency LiDAR has and a single-image model
lacks. This is the part of the oracle$\rightarrow$LiDAR gap (\SI{\rsORACLEGAPLIDARCM}{cm}) that scale alignment cannot fix, and the
reason systems like SimpleRecon use multi-image cost volumes (\S\ref{sec:related-mv}). The sparse row inherits it: its remaining gap to
LiDAR (\SI{\rsSPARSEGAPLIDAR}{cm}) is this shape term.

\subsection{Conclusions for system design}
\begin{table*}[t]
\centering\small
\begin{tabularx}{\textwidth}{@{}XX@{}}
\toprule
Assumed beforehand & What the data say \\
\midrule
mono's problem is coarse shape & shape is \SI{\rsORACLEGAPLIDARCM}{cm} behind LiDAR; scale is \SI{18}{cm} behind \\
one calibration per room suffices & no: \SI{\rsSCENECM}{cm} in every room; scale is view-dependent \\
a metric model removes the scale problem & no: but it reduces it to one parameter per frame \\
a metric model cannot be fixed without a laser scanner & largely can: a phone-LiDAR teacher takes it from \SI{\rsMETRICCM}{cm} to \SI{\rsFTLIDARALLCM}{cm} in every room (\S\ref{sec:main}); the rest is per-frame scale (\S\ref{sec:whyft}) \\
a few sparse metric points cannot fix a dense map & they can: 200 points/frame $\rightarrow$ within \SI{\rsSPARSEGAP}{cm} of the oracle \\
a bigger model means a more accurate room & $+$\SI{1.0}{cm} S$\rightarrow$L vs.\ $+$\SI{17.6}{cm} from scale \\
more frames mean more accuracy & more frames mean more coverage; mono accuracy improves with fewer frames \\
LiDAR confidence maps are free accuracy & no measurable gain from masking or weighting (Exp.~5) \\
LiDAR's hard scenes are mono's hard scenes & no (blank bedroom: hard for mono, easy for LiDAR) \\
\bottomrule
\end{tabularx}
\end{table*}
